% TOP_PROBLEM: {"id":30007578,"problem_number":"LOCAL-30007578","title":"Inequality-aware allocation","category_id":21,"category":{"id":21,"name":"economics","display_name":"Economics","description":"Open economic theory statements, published research agendas, and proposed scoped specifications; question provenance and historical priority are qualified per record.","slug":"economics","order_index":21,"created_at":"2026-10-08T00:00:00Z"},"status":"research_question","external_url":null,"legacy_ids":[],"published":false,"statement":"Characterize redistributive allocation mechanisms robust to moment-constrained uncertainty about willingness to pay and welfare weights, including the value of rationing over posted prices.","economics":{"original_id":67,"field":"Mechanism design and market design","kind":"design","formalization_basis":"adapted_research_question","source_status":"research_agenda","priority_status":"earliest_found","priority_note":"This is the earliest source in which the relevant agenda was verified during this audit; absolute priority has not been established.","earliest_verified_reference":{"authors":["Piotr Dworczak"],"title":"Inequality and Market Design","year":2024,"url":"https://www.sigecom.hosting.acm.org/exchanges/volume_22/1/issue.pdf","doi":null,"locator":"§4 Directions for Future Research, pp.89–90, especially items2–4","evidence_summary":"The letter explicitly asks about simple-mechanism guarantees, moment-based robustness, and multidimensional screening with redistributive objectives."},"current_reference":{"authors":["Piotr Dworczak"],"title":"Inequality and Market Design","year":2024,"url":"https://www.sigecom.hosting.acm.org/exchanges/volume_22/1/issue.pdf","doi":null,"locator":"§4 Directions for Future Research, pp.89–90, especially items2–4","evidence_summary":"The letter explicitly asks about simple-mechanism guarantees, moment-based robustness, and multidimensional screening with redistributive objectives."},"formalization_note":"Adapted from §4 item3 on moment-based robustness, with item2 simple-mechanism comparisons. Year is2024, correcting the earlier2025 label."},"statusQualification":"Published question or research agenda verified at the cited locator. The mathematical specification is adapted for this catalogue and includes additional assumptions; source publication does not establish first-ever priority or certify this exact model remains unsolved. Adapted from §4 item3 on moment-based robustness, with item2 simple-mechanism comparisons. Year is2024, correcting the earlier2025 label.","statusReviewedAt":"2026-10-08"}
% FIRST_POSTED: 2026-10-08
% ENTRY_KIND: statement_only
% !TeX program = lualatex
\documentclass[11pt]{article}
\usepackage{fontspec}
\usepackage[a4paper,margin=25mm]{geometry}
\usepackage{amsmath,amssymb,mathtools}
\usepackage{xurl}
\usepackage[hidelinks]{hyperref}
\newcommand{\sref}[2]{\hyperlink{source:#1:#2}{[S#2]}}
\setlength{\emergencystretch}{3em}
\begin{document}
% problemId:30007578
\section*{Inequality-aware allocation}\label{30007578}
\subsection{Definitions and mathematical statement}
A unit mass of agents demand at most one homogeneous indivisible good. Type \((r,w)\) comprises private willingness to pay \(r\in[0,1]\) and social welfare weight \(w\in[0,W]\), where finite \(W>0\) is known. The designer has supply \(q\in(0,1)\) and knows only specified moments of the joint type distribution; let \(\mathcal F\) be all distributions consistent with them. A direct mechanism specifies allocation probability \(x(r)\in[0,1]\) and payment \(p(r)\), with utility \(r x(r)-p(r)\), truthful reporting, voluntary participation, and feasibility under every \(F\in\mathcal F\). With shadow value \(\lambda\geq0\) for public revenue, evaluate
\[
 V_F(M)=E_F[w(r x(r)-p(r))+\lambda p(r)],\qquad
 \sup_M\inf_{F\in\mathcal F}V_F(M).
\]
Characterize robust optimal mechanisms and when price menus, subsidies, or rationing improve this redistributive objective. Derive sharp guarantees for posted-price mechanisms relative to the full robust optimum, under clearly stated budget constraints and welfare-weight bounds. Explain why privately varying \(w\) affects the designer's objective without being directly reportable in utility; additional screening dimensions require an extended model. Existing Bayesian one-dimensional solutions are maintained. The cited letter explicitly asks how unknown joint distributions and multidimensional extensions change inequality-aware design; the displayed moment-based robust problem makes one of those directions precise.

\par\textbf{Formulation and scope.} Adapted from §4 item3 on moment-based robustness, with item2 simple-mechanism comparisons. Year is2024, correcting the earlier2025 label.

\par\textbf{Open-question provenance.} An explicit question or research agenda was verified at the source locator. This does not by itself certify that every later version remains unresolved \sref{30007578}{1}.

\par\textbf{Historical priority.} This is the earliest source in which the relevant agenda was verified during this audit; absolute priority has not been established.

\subsection{Short English statement}
Characterize redistributive allocation mechanisms robust to moment-constrained uncertainty about willingness to pay and welfare weights, including the value of rationing over posted prices.

\subsection{Sources}
\begin{itemize}\raggedright
\item[\textnormal{[S1]}]\hypertarget{source:30007578:1}{} Piotr Dworczak. 2024. Inequality and Market Design. §4 Directions for Future Research, pp.89–90, especially items2–4.
\url{https://www.sigecom.hosting.acm.org/exchanges/volume_22/1/issue.pdf}.
{\small\textit{Earliest verified question/agenda reference; historical priority qualified below. The letter explicitly asks about simple-mechanism guarantees, moment-based robustness, and multidimensional screening with redistributive objectives.}}
\end{itemize}
\end{document}
